\section{Conclusion}
\label{sec:conclusion}

We examine Indonesia's rupiah cryptocurrency premium using a fixing-synchronized panel from January 2021 through November 2026 ($n = 1,414$ Bank Indonesia JISDOR business days across BTC, ETH, and XRP, paired with $616,683$ regular five-minute bars matched from $616,896$ raw Indodax and $616,683$ raw Binance bars) alongside reconstructed historical tax and exchange-fee schedules. Pre-November 2025 historical series rest on Indodax alone, complemented by 349 days of direct Binance IDR cross-checks (November 29, 2025--November 12, 2026) and a single-day November 13, 2026 order-book depth snapshot.

Four empirical findings stand out. First, once global \texttt{BTC/USDT} is converted through Indodax's synchronous \texttt{USDT/IDR} book, the pure cross-exchange crypto leg averages $-0.83\text{ bps}$ (statistically indistinguishable from zero), and $98.44\%$ of business-day deviations sit inside reconstructed dynamic domestic two-leg taker cost bands (sensitivity range $[97.52\%, 98.51\%]$ across partially verified 2022--2024 fee tiers; $97.81\%$ against the fixed $66.45\text{ bps}$ sample-average taker threshold). Second, by accounting identity, the domestic \texttt{USDT/IDR} minus JISDOR spread accounts for $83.10\%$ of total BTC basis variance ($+30.77\text{ bps}$ mean, $73.66\text{ bps}$ SD), while pure crypto deviations account for $11.53\%$ and their covariance $5.37\%$. Volume-averaging \texttt{USDT/IDR} over each day's exact Bank Indonesia JISDOR monitoring window leaves the stablecoin--FX share at $82.39\%$, so the baseline share is not an averaging-window artifact.

Third, separating Indodax \texttt{BTC/IDR}, Binance \texttt{BTC/USDT}, and Indodax \texttt{USDT/IDR} in a three-variable five-minute VECM on the regular clock-time grid ($T_{eff} = 616,677$ bars) shows that we fail to reject weak exogeneity for global Binance \texttt{BTC/USDT} ($\hat{\alpha} = +0.0011$, HAC $p = 0.418$; Hasbrouck IS $[42.6\%, 99.0\%]$), whereas Indodax \texttt{BTC/IDR} bears $80.0\%$ ($\hat{\alpha} = -0.0445$) and Indodax \texttt{USDT/IDR} bears $18.0\%$ ($\hat{\alpha} = +0.0100$) of short-run error correction, with a full-system companion half-life of $25.08$ five-minute bars (scalar $1+\beta'\alpha$ benchmark $12.09$ bars). Fourth, under controlled $\pm 30$-day HAC regressions with Holm adjustment across the 20-test event family, only the November 1, 2024 CFX/KKI fee launch survives at $5\%$ on $B_t^{stablecoin\text{-}FX}$ ($+17.37\text{ bps}$, $p_{\text{Holm}} = 0.026$, a controlled post-event association rather than a one-for-one fee pass-through), while PMK 50/2025's shift on $B_t^{crypto}$ is specification-sensitive ($+21.23\text{ bps}$, $p_{\text{Holm}} = 0.011$ unadjusted vs.\ $+17.22\text{ bps}$, $p_{\text{Holm}} = 0.275$ controlled). Measured at the JISDOR fixing, Indonesia's headline bitcoin premium is overwhelmingly a domestic stablecoin--FX spread rather than a cross-venue bitcoin price gap.
